% Kapitel 3 — Budget: 8 Seiten — Schreiben: 14.–23.12.
% Text liegt vor (Exposé), eher kürzen als ergänzen; nur, was im Moment für diese
% Arbeit relevant ist (Beschluss 26.10.). Stand vom Anfang der Arbeit, wird nicht
% nachgezogen; die Schwenks (Ausrichtungsänderungen) gehören NICHT hierher.

\chapter{Stand der Forschung}
\label{chap:forschungsstand}

\section{Kontext als Erfolgsfaktor und Engpass}% 1,5 S., bisher 2.2
\label{sec:kontext-engpass}

\section{Graphbasierte Kontextbereitstellung}% 2 S., bisher 2.3
\label{sec:graphbasiert}

\section{Grenzen langer, unstrukturierter Kontexte (Modellseite)}% 1 S., bisher 2.4, kürzen
\label{sec:grenzen-kontexte}

\section{Nachvollziehbarkeit und menschliche Beteiligung}% 1 S., bisher 2.5, kürzen
\label{sec:nachvollziehbarkeit}

\section{Abgrenzung und Beitrag}% 2,5 S., bisher Kap. 3 (Beschluss 26.10.)
% Inhalt: automatisches graphbasiertes Code-Retrieval, interaktive Inspektion von
% Retrieval-Pipelines, menschengestützte Interfaces für Wissensgraphen.
\label{sec:abgrenzung}
